% !TEX program = xelatex
% 공통수학2 · Ⅰ. 도형의 방정식 · 2. 원의 방정식 · 02 원과 직선의 위치 관계 · 1/2차시
\documentclass[11pt,a4paper]{article}
\usepackage{kotex}
\usepackage{iftex}
\ifPDFTeX\else
  \setmainhangulfont{Noto Serif CJK KR}
  \setsanshangulfont{Noto Sans CJK KR}
\fi
\usepackage[margin=2.1cm]{geometry}
\usepackage{parskip}
\usepackage{enumitem}
\usepackage{array}
\usepackage{tabularx}
\usepackage{booktabs}
\usepackage{xcolor}
\usepackage{amssymb}
\usepackage{tikz}
\usepackage[most]{tcolorbox}
\usepackage{fancyhdr}
\usepackage{titlesec}

\definecolor{goldc}{HTML}{B07C22}
\definecolor{goldstrong}{HTML}{8A5F16}
\definecolor{indigoc}{HTML}{2B3B57}
\definecolor{indigosoft}{HTML}{6E82A6}
\definecolor{linec}{HTML}{D6DACB}
\definecolor{surface2}{HTML}{E4E8DC}
\definecolor{notebg}{HTML}{FBF3E3}
\definecolor{inksoft}{HTML}{52604F}

\titleformat{\section}{\normalfont\sffamily\bfseries\large\color{indigoc}}{\thesection}{0.6em}{}
\titlespacing{\section}{0pt}{16pt}{8pt}
\newtcolorbox{ideabox}{enhanced, breakable, colback=white, colframe=goldc, boxrule=0.5pt, leftrule=3pt, arc=2pt, left=10pt, right=10pt, top=8pt, bottom=8pt}
\newtcolorbox{calloutbox}{enhanced, breakable, colback=white, colframe=indigosoft, boxrule=0.5pt, leftrule=3pt, arc=2pt, left=10pt, right=10pt, top=7pt, bottom=7pt}
\newtcolorbox{notebox}{enhanced, breakable, colback=notebg, colframe=goldc, boxrule=0.5pt, leftrule=3pt, arc=2pt, left=10pt, right=10pt, top=7pt, bottom=7pt}
\newtcolorbox{answerbox}{enhanced, breakable, colback=white, colframe=linec, boxrule=0.5pt, arc=2pt, left=10pt, right=10pt, top=7pt, bottom=7pt}
\newcommand{\lessonbanner}[3]{%
  \begin{tcolorbox}[enhanced, colback=indigoc, colframe=indigoc, boxrule=0pt, arc=0pt,
    left=14pt, right=14pt, top=14pt, bottom=10pt, borderline south={2.4pt}{0pt}{goldc}, fontupper=\color{white}]
    {\sffamily\footnotesize\color{white!85!black} #1}\\[4pt]{\sffamily\bfseries\Large #2}\\[3pt]{\sffamily\small\color{white!88!black} #3}
  \end{tcolorbox}}
\pagestyle{fancy}\fancyhf{}\renewcommand{\headrulewidth}{0pt}
\fancyfoot[L]{\footnotesize\color{inksoft} 공통수학2 · 2026학년도 2학기 · 지평선고등학교}
\fancyfoot[R]{\footnotesize\color{inksoft} 교사용 지도안 -- 배포 금지}

\begin{document}
\lessonbanner{공통수학2 · Ⅰ. 도형의 방정식 · 2. 원의 방정식}{02 원과 직선의 위치 관계 --- 1/2차시}{판별식으로 위치 관계 판단하기 · 교사용 지도안 (2026학년도 2학기)}
\vspace{10pt}

\noindent\begin{tabularx}{\linewidth}{@{}>{\bfseries\color{indigoc}}p{2.6cm}X@{}}
\toprule
대상 & 고1 · 2개 반 · 반당 13명 \\
차시 & 50분 (도입 10 · 전개 30 · 정리 10) \\
성취기준 & 좌표평면에서 원과 직선의 위치 관계를 판단하고, 이를 활용하여 문제를 해결할 수 있다. \\
선행 학습 & 9$\sim$10차시 --- 원의 방정식 \\
\bottomrule
\end{tabularx}

\section{학습 목표 \& Big Idea}
\begin{ideabox}
\textbf{\color{goldstrong}영속적 이해 (Big Idea)}\\
원과 직선의 `만남'이라는 기하적 사건은 연립방정식의 해의 개수, 즉 이차방정식의 판별식이라는 대수적 질문으로 정확히 번역된다.
\vspace{4pt}
\small\color{inksoft} 핵심개념: \textbf{판별식} \quad\textbullet\quad 핵심개념: \textbf{두 가지 방법} \quad\textbullet\quad 일반화: \textbf{같은 질문을 여러 도구로 풀 수 있다}
\end{ideabox}
\begin{calloutbox}
\textbf{차시 학습 목표}
\begin{itemize}[leftmargin=1.4em, itemsep=2pt, topsep=4pt]
  \item 원과 직선의 방정식을 연립하여 얻은 이차방정식의 판별식으로 위치 관계를 판단할 수 있다.
  \item 원의 중심과 직선 사이의 거리($d$)와 반지름($r$)의 대소 비교로도 같은 판단을 할 수 있음을 안다.
\end{itemize}
\end{calloutbox}

\section{질문의 연결}
\begin{tabularx}{\linewidth}{@{}>{\bfseries\color{indigoc}\small}p{2.4cm}X@{}}
사실적 질문 & 원과 직선이 만나는 점의 개수는 어떻게 판단할 수 있을까? \\[6pt]
개념적 질문 & 판별식 부호(대수)와 $d$, $r$의 대소(기하), 두 방법은 왜 항상 같은 결론에 도달할까? \\[6pt]
논쟁적 질문 & 별똥별이 보름달을 스치듯 지나가는 순간처럼, `접한다'는 것은 `만난다'와 `만나지 않는다' 그 경계에서 어떤 의미를 가질까? \\
\end{tabularx}

\section{수업 흐름}
\begin{tabularx}{\linewidth}{@{}>{\centering\arraybackslash}p{1.6cm}X>{\raggedright\arraybackslash\small}p{3.1cm}@{}}
\toprule
\textbf{단계} & \textbf{활동 \& 발문} & \textbf{ATL / VTR} \\
\midrule
\textcolor{goldstrong}{\textbf{도입}}\newline\footnotesize 10분 &
\textbf{마음공부 (2분)} --- 유무념 대조.\newline
\textbf{생각 열기 (8분)} --- 반지름 $r$인 원과 세 직선(각각 원의 중심에서 거리 $d$가 다름)을 제시 $\to$ $d$와 $r$의 대소에 따라 $>,=,<$ 중 알맞은 것을 채우며 직관적으로 위치 관계를 판단. &
ATL: 사고(패턴 발견)\newline VTR: See-Think-Wonder \\
\addlinespace
\textcolor{goldstrong}{\textbf{전개}}\newline\footnotesize 30분 &
\textbf{개념 형성 (15분)} --- 원 $x^2+y^2=r^2$과 직선 $y=mx+n$을 연립하여 이차방정식 $(m^2+1)x^2+2mnx+n^2-r^2=0$을 얻고, 판별식 $D$의 부호로 교점의 개수가 정해짐을 유도 $\to$ 예제 1을 두 가지 방법(판별식, 거리 비교)으로 함께 풀이.\newline
\textbf{적용 (15분)} --- 예제 2(만나도록 하는 $k$의 범위) 함께 풀이 $\to$ 문제 1, 2 짝 활동(두 방법 모두 시도해 보기). &
ATT: 촉진자 역할\newline ATL: 사고·비교 \\
\addlinespace
\textcolor{goldstrong}{\textbf{정리}}\newline\footnotesize 10분 &
\textbf{개념 연결 질문 (5분)} --- ``판별식 방법과 거리 비교 방법, 각각 언제 더 편리할까?''\newline
\textbf{성찰 (5분)} --- 감사 일기 1문장 + 다음 차시 예고(원의 접선의 방정식). &
마음공부 · 감사 일기 \\
\bottomrule
\end{tabularx}

\begin{calloutbox}
\textbf{지도상의 유의점} --- 두 방법(판별식, 거리 비교)이 결국 같은 결론임을 예제 1에서 나란히 보여 주어, 학생이 상황에 맞게 방법을 선택할 수 있게 한다. 접한다(=한 점에서 만난다)는 것이 $D=0$ 또는 $d=r$과 정확히 대응됨을 강조한다.
\end{calloutbox}

\section{형성평가 \& 답안}
\begin{answerbox}
\textbf{예제 2} \quad 원 $x^2+y^2=9$와 직선 $y=x+k$가 서로 다른 두 점에서 만나는 $k$의 범위\\
\textbf{\color{goldstrong}답} \; $-3\sqrt2 < k < 3\sqrt2$
\end{answerbox}
\begin{minipage}[t]{0.48\linewidth}
\begin{answerbox}
\textbf{문제 1}\\
\small ⑴ $x^2+y^2=1$, $y=-2x+3$ \quad ⑵ $x^2+y^2=8$, $x+y=4$\\[4pt]
\textbf{\color{goldstrong}답} \; ⑴ 만나지 않는다 \quad ⑵ 접한다(한 점에서 만난다)
\end{answerbox}
\end{minipage}\hfill
\begin{minipage}[t]{0.48\linewidth}
\begin{answerbox}
\textbf{문제 2} \quad 원 $x^2+y^2=4$와 직선 $y=2x+k$가 만나지 않는 $k$의 범위\\[4pt]
\textbf{\color{goldstrong}답} \; $k<-2\sqrt5$ 또는 $k>2\sqrt5$
\end{answerbox}
\end{minipage}

\section{성취수준 (이 차시 범위)}
\begin{tabularx}{\linewidth}{@{}>{\bfseries\color{goldstrong}}p{0.9cm}X@{}}
\toprule
A & 원과 직선의 위치 관계를 다양한 방법으로 판단하고, 문제를 해결할 수 있다. \\
B & 원과 직선의 위치 관계를 다양한 방법으로 판단하고, 간단한 문제를 해결할 수 있다. \\
C & 원과 직선의 위치 관계를 다양한 방법으로 판단할 수 있다. \\
D & 판별식을 이용해 원과 직선의 위치 관계를 판단할 수 있다. \\
E & 원과 직선의 위치 관계를 안다. \\
\bottomrule
\end{tabularx}

\section{다음 차시}
\begin{calloutbox}
\textbf{02 원과 직선의 위치 관계 --- 2/2차시.} 원의 접선의 방정식(기울기가 주어진 경우, 원 위의 한 점에서의 경우, 원 밖의 한 점에서 그은 경우)을 다룬다.
\end{calloutbox}
\end{document}
